\chapter{Documentation}
\label{chp:documentation} 

This is the second of three chapters that aim to move from writing code
that works to writing software. In this particular chapter you will consider how
to write documentation for your code.

\begin{note}
In this chapter you will cover:
\begin{itemize}
\item 

Using the \index{Diataxis} framework for
documentation~\cite{Procida_Diataxis_documentation_framework}.

\end{itemize}
\end{note}




\section{Tutorial}
\label{sec:documentation_tutorial}

In this tutorial you will write documentation for the code
you wrote in Section~\ref{sec:modularisation_tutorial}.


You start by creating a new file in VScode called \texttt{README.md}.

\begin{note}
You will be writing your documentation in markdown.
\end{note}


Start by writing the title of your library and a quick single sentence
description.

\begin{md}
# Absorption

Functionality to study the absorbing Markov chains.
\end{md}


\subsection{Writing a tutorial}
\label{\detokenize{building-tools/06-documentation/tutorial/main:writing-a-tutorial}}

You will then write your first section which is a \textbf{tutorial}.

\begin{note}
The goal of a tutorial is to provide a hands on introduction and demonstration
of the software.
\end{note}

\begin{md}

In this tutorial we will see how to use `absorption` to study an absorbing
Markov chain. For some background information on absorbing Markov chains we
recommend: <https://en.wikipedia.org/wiki/Absorbing_Markov_chain>.

Given a transition matrix $P$ defined by:

$$
p = \begin{pmatrix}
   1/2 & 1/4 & 1/4\\
   1/3 & 1/3 & 1/3\\
   0   & 0   & 1
    \end{pmatrix}
$$

We will start by seeing how the chain evolves over time by starting with an
initial vector $\pi=(1,0,0)$. In the next code snippet we will import the
necessary libraries and create both $P$ and $\pi$:

```python
import numpy as np

import absorption

pi = np.array([1, 0, 0])
P = np.array([[1 / 2, 1 / 4, 1 / 4], [1 / 3, 1 / 3, 1 / 3], [0, 0, 1]])
```

We now see how the vector $\pi$ changes over time:

```python
for k in range(10):
    print(absorption.get_long_run_state(pi, k, P))
```

This will give:

```
[1. 0. 0.]
[0.5  0.25 0.25]
[0.33333333 0.20833333 0.45833333]
[0.23611111 0.15277778 0.61111111]
[0.16898148 0.1099537  0.72106481]
[0.12114198 0.0788966  0.79996142]
[0.08686986 0.05658436 0.85654578]
[0.06229638 0.04057892 0.8971247 ]
[0.0446745 0.0291004 0.9262251]
[0.03203738 0.02086876 0.94709386]
```

We see that, as expected over time the probability of being in the third state,
which is absorbing, increases.

We can also use `absorption` to get the average number of steps until
absorption from each state:

```python
absorption.compute_t(P)
```

This gives:

```
array([3.66666667, 3.33333333])
```

We see that the expected amounts of steps from the first state is slightly more
than from the second.
\end{md}

This \textbf{tutorial} section allows newcomers to your code to see how it is intended
to be used.

\subsection{Writing the how-to guides}
\label{\detokenize{building-tools/06-documentation/tutorial/main:writing-the-how-to-guides}}

In the next section you will write a series of \textbf{how to} guides, this is
targeted
at someone who has perhaps worked through the tutorial already and wants to
directly know how to do a specific task.

Directly underneath what you have written so far write:

\begin{md}

## How to guides

### How to compute the long run state of a system after a given number of steps

Given a transition matrix $P$ and a state vector $\pi$, the state of the system
after $k$ steps is given by:

```python
import numpy as np

import absorption

pi = np.array([0, 1, 0])
P = np.array([[1 / 3, 1 / 3, 1 / 3], [1 / 4, 1 / 4, 1 / 2], [0, 0, 1]])
absorption.get_long_run_state(pi=pi, k=10, P=P)
```

This gives:

```
array([0.0019552, 0.0019552, 0.9960896])
```

### How to extract the transitive state transition sub matrix $Q$

Given a transition matrix $P$, the sub matrix $Q$ that
corresponds to the transitions between transitive (i.e. not absorbing) states can
be extracted:

```python
import numpy as np

import absorption

P = np.array([[1 / 3, 1 / 3, 1 / 3], [0, 1, 0], [1 / 4, 1 / 4, 1 / 2]])
absorption.extract_Q(P=P)
```

This gives:

```
array([[0.33333333, 0.33333333],
       [0.25      , 0.5       ]])
```

### How to compute the fundamental matrix $N$

Given a transition matrix $P$, the fundamental matrix $N$
can be obtained:

```python
import numpy as np

import absorption

P = np.array([[1 / 3, 1 / 3, 1 / 3], [0, 1, 0], [1 / 4, 1 / 4, 1 / 2]])
Q = absorption.extract_Q(P=P)
absorption.compute_N(Q=Q)
```

This gives:

```
array([[2.        , 1.33333333],
       [1.        , 2.66666667]])
```

### How to compute the average steps until absorption

Given a transition matrix $P$ and a state vector $\pi$, the average number of
steps until absorption from all states can be obtained:

```python
import numpy as np

import absorption

P = np.array([[1 / 3, 1 / 3, 1 / 3], [0, 1, 0], [1 / 4, 1 / 4, 1 / 2]])
absorption.compute_t(P=P)
```

This gives:

```
array([3.33333333, 3.66666667])
```
\end{md}

This \textbf{how to} section is an efficient collection of instructions to be
able to carry out specific tasks made possible by the software.


\subsection{Writing the explanations section}
\label{\detokenize{building-tools/06-documentation/tutorial/main:writing-the-explanations-section}}

In the next section you will write the \textbf{explanations} which aims to give more
in depth understanding not necessarily directly related to the code.

\begin{md}
## Explanation

### Brief overview of absorbing markov chains

A Markov chain with a given transition matrix $P$ is a system that moves from
state to state randomly with the probabilities given by $P$.

For example:

$$
P = \begin{pmatrix}
        1 / 3 & 1 / 3 & 1 / 3 \\
        0     & 1     & 0     \\
        1 / 4 & 1 / 4 & 1 / 2
    \end{pmatrix}
$$

The corresponding Markov chain has 3 states and:

- $P_{11}=1/3$ means that when in state 1 there is a 1/3 chance of staying in
  state 1.
- $P_{23}=0$ means that when in state 2 there is a 0 chance of staying in
  state 1.
- $P_{22}=$ actually implies that once we are in state 2 that the only next
  state is state 2.

In general:

$$
    P_{ij} > 0 \text{ for all }ij
$$

$$
    \sum_{j=0}^{|P|} P_{ij} = 1 \text{ for all }i
$$

If $P_{ii}=1$ then state $i$ is an absorbing state from which no further changes
can occur.

In the case of absorbing markov chains there are a number of quantities that can
be measured.

### Calculating the state after a given number of iterations

Given a vector that describes the state of the system $\pi$ and a transition
matrix $P$, the state of the system after $k$ iterations will be given by the
following vector:

$$
    \pi P ^ k
$$

### The canonical form of the transition matrix

A transition matrix $P$ is written in its canonical form when it can be written
as:

$$
  P =
  \left(\begin{array}{c|c}
    Q & R \\\hline
    0 & I
  \end{array}\right)
$$

Where $Q$ is the matrix of transitions between non absorbing states.

For example, the canonical form of $P$ would be:

$$
    \begin{pmatrix}
        1 / 3 & 1 / 3 & 1 / 3 \\
        1 / 4 & 1 / 2 & 1 / 4 \\
        0     &  0    & 1     \\
    \end{pmatrix}
$$

which would give:

$$
Q = \begin{pmatrix}
    1 / 3 & 1 / 3 \\
    1 / 4 & 1 / 2
    \end{pmatrix}
$$

### The fundamental matrix

Given $Q$ the fundamental matrix $N$ is defined as:

$$N = (I - Q) ^{-1}$$

$N_{ij}$ corresponds to the expected number of times the chain will be in state
$j$ given that it started in state $i$.

### The expected number of steps until absorption.

Given $N$, the expected number of steps until absorption is given by the vector:

$$
t = N \mathbb{1}
$$

where $\mathbb{1}$ denotes a vector of 1s.
\end{md}

This \textbf{explanations} section gives background reading as to how the code works.


\subsection{Writing the reference section}
\label{\detokenize{building-tools/06-documentation/tutorial/main:writing-the-reference-section}}

In the next section you will write the \textbf{reference} which aims to be a concise
collection of reference material.

\begin{md}
## Reference

### List of functionality

The following functions are written in `absorption`:

- `get_long_run_state`
- `extract_Q`
- `compute_N`
- `compute_t`

### Bibliography

The wikipedia page on absorbing Markov chains gives a good overview of the
topic: <https://en.wikipedia.org/wiki/Absorbing_Markov_chain>

The following text is a recommended reference on Markov chains:

> Stewart, William J. Probability, Markov chains, queues, and simulation: the
> mathematical basis of performance modelling. Princeton university press, 2009.
\end{md}

Figure~\ref{fig:documentation_in_vscode} shows the start of the markdown file
in VScode alongside the preview. Note that the \texttt{Markdown all in one} plugin ensures that
the mathematics is rendered see Section~\ref{sec:plugins_for_vs_code} for
information on installing plugins.

\begin{figure}[htbp]
\centering


\includegraphics[width=0.750\linewidth]{./assets/documentation_in_vscode/main.png}
\caption{The \texttt{README.md} file in VScode with the rendered preview (using
the \texttt{Markdown all in one} plugin).}
\label{fig:documentation_in_vscode}
\end{figure}

\section{How to}
\label{\detokenize{building-tools/06-documentation/how/main:how}}\label{\detokenize{building-tools/06-documentation/how/main::doc}}

\subsection{Write documentation}
\label{\detokenize{building-tools/06-documentation/how/main:how-to-write-documentation}}

Follow the Diataxis framework~\cite{Procida_Diataxis_documentation_framework} for documentation.

This involves separating your documentation in to 4 different sections based on
separate aims for readers.
\begin{itemize}
\item 

Tutorial: for learning.

\item 

How to guides: to achieve a specific goal.

\item 

Explanation: to understand.

\item 

Reference: to find information.

\end{itemize}


\subsubsection{Write a tutorial}

A tutorial should include step by step instructions with expected behaviours.
This should not focus on any deeper explanation.

An analogy of this is teaching a young toddler to build a toy train
track. They do not need to know the physics related to how the train will go
through the track. They need only to see how to lay the track pieces.

\subsubsection{How to write a how to guide}
\label{\detokenize{building-tools/06-documentation/how/main:how-to-write-a-how-to-guide}}\label{\detokenize{building-tools/06-documentation/how/main:id2}}

A how to guide should provide a quick and to the point description of how to
solve a specific problem.

An analogy of this would be a recipe. The recipe will not necessarily explain
how to chop an onion and/or why you are chopping an onion. It will tell you
how to chop an onion as a step of cooking a particular meal.

\subsubsection{How to write an explanation section}
\label{\detokenize{building-tools/06-documentation/how/main:how-to-write-an-explanation-section}}\label{\detokenize{building-tools/06-documentation/how/main:id3}}

The explanation section should provide a deeper understanding of the concepts
under the code.

An analogy of this again related to a recipe would be a book on the chemistry of
taste and why a chopped onion adds a specific type of flavour to a meal.

\subsubsection{How to write a reference section}
\label{\detokenize{building-tools/06-documentation/how/main:how-to-write-a-reference-section}}\label{\detokenize{building-tools/06-documentation/how/main:id4}}

The reference section should provide an overview of the specific tools, commands
and indeed place for background reading as well (although this can also be
referred to in the explanation section).

\subsection{Write markdown}
\label{sec:how_to_write_markdown}

\subsubsection{How to include section headers in markdown}
\label{\detokenize{building-tools/06-documentation/how/main:how-to-include-section-headers-in-markdown}}\label{\detokenize{building-tools/06-documentation/how/main:id5}}

To include a section header in markdown use \texttt{\#}. The number of \texttt{\#} corresponds
to the level of the section header.

\begin{md}
# Section
\end{md}


For example:

\begin{md}
# The absorption library

Functionality to study the absorbing Markov chains.
\end{md}




\subsubsection{Include code in markdown}
\label{\detokenize{building-tools/06-documentation/how/main:how-to-include-code-in-markdown}}\label{\detokenize{building-tools/06-documentation/how/main:id6}}

To include code in markdown use three ` marks followed by the name of the
language:

\begin{md}
```<language>

<code>
```
\end{md}


For example:

\begin{md}
```python
import sympy as sym

x = sym.Symbol(''x'')
```
\end{md}


It is also possible to include code in markdown using an indented block:


For example:

\begin{md}
Here is some code:

    import sympy as sym

    x = sym.Symbol(''x'')
\end{md}


\begin{note}
Using an indented block does not allow you to specify the language and can lead
to mistake when combining with other nested statement.
\end{note}


\subsubsection{Include a hyperlink in markdown}
\label{\detokenize{building-tools/06-documentation/how/main:how-to-include-a-hyperlink-in-markdown}}\label{\detokenize{building-tools/06-documentation/how/main:id7}}

To include a hyperlink in markdown use \texttt{{[}{]}()}
language:

\begin{md}
[text](url)
\end{md}


For example:

\begin{md}
The [Online Encyclopedia of Integer Sequences](https://oeis.org) is a good resources for studying
resources.
\end{md}


\subsubsection{Include an image in markdown}
\label{\detokenize{building-tools/06-documentation/how/main:how-to-include-an-image-in-markdown}}\label{\detokenize{building-tools/06-documentation/how/main:id8}}

To include an image in markdown use \texttt{!{[}{]}()}:

\begin{md}
![caption](path)
\end{md}


For example:

\begin{md}
Here is an image:

![An image](image.jpg)
\end{md}

\begin{note}

If the image file is not located in the same directory as the markdown file the
path to the file must be correct.
\end{note}




\section{Exercises}

Write documentation for the \texttt{statistics.py} file written in the exercises of
Chapter~\ref{chp:modularisation}.


\section{Further information}
\label{\detokenize{building-tools/06-documentation/why/main:further-information}}\label{\detokenize{building-tools/06-documentation/why/main:documentation-further-information}}\label{\detokenize{building-tools/06-documentation/why/main::doc}}

\subsection{What is documentation?}
\label{\detokenize{building-tools/06-documentation/why/main:what-is-documentation}}

Documentation can have many different interpretations. A good definition is
given in \cite{martraire2019living}.

\begin{quote}
``The process of transferring valuable knowledge to other people now and also to
people in the future.''
\end{quote}

\begin{note}
It is important to realise that the target of the documentation can be the
writer of the software itself at a future date.
\end{note}


There are two types of documentation:
\begin{itemize}
\item 

\textbf{Internal documentation} which includes things like docstrings and a good
choice of variable names.

\item 

\textbf{External documentation} which includes things like \texttt{README.md} and other
separate documentation.

\end{itemize}


For a software project to be well documented it needs \textbf{both} internal and
external documentation.


In~\cite{martraire2019living} there are 
4 properties of documentation:
\begin{itemize}
\item 

Reliable: it needs to be accurate.

\item 

Low effort: it should require minimal effort when changes are made to the code
base.

\item 

Collaborative: it should be a tool from which collaboration can occur.

\item 

Insightful: it should give information not only to be able to use the code but
also to understand specific reasons why certain decisions have been made as to
its design.

\end{itemize}


\subsection{What is the purpose of the four separate sections in documentation?}

As discussed in~\cite{Procida_Diataxis_documentation_framework}:
\begin{quote}


``Tutorials are lessons that take the reader by the hand through a series of
steps to complete a project of some kind. They are what your project needs in
order to show a beginner that they can achieve something with it.''
\end{quote}


\begin{quote}
``How-to guides take the reader through the steps required to solve a
real-world problem'''
\end{quote}

\begin{quote}
``Reference guides are technical descriptions of the machinery and how to
operate it.''
\end{quote}

\begin{quote}
``Explanation, or discussions, clarify and illuminate a particular topic. They
broaden the documentation’s coverage of a topic.''
\end{quote}


It is natural when describing a project for the boundaries between these four
topics to become fuzzy. Thus, having them explicitly in four separate sections
ensures the reader is able to specifically find what they need.

\subsection{What alternatives are there to writing documentation in the \texttt{README.md} file?}

A single \texttt{README.md} file is a good way to start documenting code. However as a
project grows it could be beneficial to use some other tools. One such example
of this is to use \texttt{sphinx}. This uses a
different markup language called \textbf{restructured text}
 and helps build more complex
documents but also interfaces to the code itself if necessary. For example, it
is possible to include the code docstrings directly in the documentation (a good
way of adding to the reference section).
